再进一步，均匀层大偏差与尺寸偏置倾斜应共同决定 oracle 容量缺口的宏观相图。

\begin{conjecture}[oracle 容量缺口指数的双 Legendre 最小包络]\label{conj:xi-time-part57eb-oracle-gap-min-envelope-phase-diagram}
沿用定理 \ref{thm:xi-time-part57eb-uniform-logmultiplicity-ldp} 与定理 \ref{thm:xi-time-part57eb-sizebias-tilt-rate-shift} 的记号。取
\[
B_m:=\lfloor \beta m\rfloor,
\qquad
\tau:=\beta\log 2,
\]
并定义上尾速率函数
\[
J_u(\tau):=\inf_{y>\tau}I_u(y),
\qquad
J_\mu(\tau):=\inf_{y>\tau}I_\mu(y).
\]
则容量缺口
\[
\mathcal G_m(B_m)=1-\mathrm{Succ}_m(B_m)
\]
的指数阶应满足
\[
\lim_{m\to\infty}-\frac1m\log \mathcal G_m(B_m)
=
\min\Bigl\{
J_\mu(\tau),\,
J_u(\tau)-\bigl(\log\varphi-\log 2+\tau\bigr)
\Bigr\}.
\]
特别地，两项相等的曲线应给出 oracle 容量缺口宏观相图的一阶相变边界。

事实上，若记
\[
\tau_m:=\frac{B_m\log 2}{m},
\]
则有精确恒等分解
\[
\mathcal G_m(B_m)
=
\mu_m(Y_m>\tau_m)
-\exp\!\Bigl(m\Bigl(\tau_m+\frac1m\log|X_m|-\log 2\Bigr)\Bigr)\,
u_m(Y_m>\tau_m).
\]
该猜想断言：在大偏差尺度上，上式两项不存在额外的指数级相消，因此缺口指数恰由两条尾率中衰减较慢者决定。
\end{conjecture}

\endinput
